\documentclass{situationist-preprint}

\title{Williams on the Thought Too Many}
\msid{ETH-19-b}
\msfamily{ETH-19}
\msdiscipline{Normative and Applied Ethics}
\mstype{Companion manuscript}
\mskeywords{partiality; one thought too many; Williams; motivation and justification; relational reasons}

\begin{document}
\maketitle

\begin{abstract}
Williams's remark that the man who saves his wife on the ground that it is permissible to do so has \enquote{one thought too many} is often read as a point about motivation, which impartialists can accommodate by keeping the impartial justification out of the agent's deliberation. We argue that the remark is a point about justification. The thought is too many because the reason it supplies is the wrong reason: a reason that would equally favor any person standing in a relationship with the same welfare effects. The reason that actually suffices, that she is his wife, is relational and non-transferable. This reading supports the conclusion of ETH-19 and explains why indirect impartialism does not answer Williams.
\end{abstract}

\section{Two readings}

Fried discussed a man who can save only one of two people in peril, one of them his wife, and allowed that he may save his wife \citep{Fried1970}. Williams objected to the suggestion that the man needs a justification, that in situations of this kind it is permissible to save one's wife: such a man has one thought too many, and the thought he should have is that it was his wife \citep[18]{Williams1981}.

On the \emph{motivational} reading, the complaint is that the man deliberates when he should simply act. Railton's sophisticated consequentialist answers it: impartial considerations justify the disposition to save one's wife, and the disposition operates without the justificatory thought \citep{Railton1984}. On the \emph{justificatory} reading, the complaint is that the permission does not rest on the impartial thought at all, so that a man who thinks it has the wrong reason, whether or not he thinks it at the moment of action. Only the second reading engages the claim of ETH-19-a that the permission is not grounded in an impartial justification (its Theorem~3.3).

\section{The argument}

\begin{proposition}\label{p:transfer}
A reason supplied by an impartial justification favors saving any person whose rescue falls under the justified practice; the reason that suffices for the man, that she is his wife, favors saving only the person with whom he stands in the relationship.
\end{proposition}

\begin{proof}
By Definition~2.2 of ETH-19-a, an impartial justification derives the permission from a principle that gives no fundamental weight to the agent's relationships, together with facts about the effects or adoptability of partial practices. Any person whose rescue would have the same effects falls under the same derivation. The reason that she is his wife is grounded in a relationship between him and her (Definition~2.1), and a qualitatively similar stranger does not stand in it.
\end{proof}

By Corollary~3.4 of ETH-19-a, the agent who saves his wife because she is his wife acts on a sufficient reason. By Proposition~\ref{p:transfer}, the impartial thought supplies a different reason, one that does not single out his wife. The man who has the impartial thought therefore takes himself to be saving his wife because she is someone the practice permits him to favor, and this misdescribes what he values. It is in this sense that the thought is too many: it is not merely superfluous in deliberation but a substitute for the reason that grounds the permission. The same structure appears in love, whose reasons, as EXI-16-a argues, are relational and do not transfer to qualitatively similar persons.

\begin{objection}
The indirect impartialist does not ask the man to think the impartial thought; he asks only that the thought be true. Williams's complaint then lapses.
\end{objection}

\begin{reply}
If the impartial thought is true and grounds the permission, the man's own reason is sufficient only derivatively, because it tracks what the justification certifies. Principle~P1 of ETH-19-a then implies that in circumstances where partial practices had worse effects the permission would fail, and Lemma~3.1 shows that it does not. The indirect view keeps the man's psychology intact at the cost of making his reason a proxy; it accepts the motivational reading of Williams and has no answer to the justificatory one \citep[see also][]{Wolf1992}.
\end{reply}

\section{Conclusion}

Williams's thought too many is not an excess of deliberation but the wrong reason. The reason that suffices is relational and non-transferable, and it needs no impartial certification. Partiality is limited by the duties of beneficence and rescue, and within the personal prerogative it may be exercised in favor of those close to one, as ETH-28-a assumes.

\begin{thebibliography}{9}\small
\bibitem[Fried(1970)]{Fried1970} Fried, C. (1970). \emph{An Anatomy of Values: Problems of Personal and Social Choice}. Harvard University Press.
\bibitem[Railton(1984)]{Railton1984} Railton, P. (1984). Alienation, consequentialism, and the demands of morality. \emph{Philosophy \& Public Affairs}, 13(2), 134--171.
\bibitem[Williams(1981)]{Williams1981} Williams, B. (1981). Persons, character and morality. In \emph{Moral Luck: Philosophical Papers 1973--1980}, Cambridge University Press, 1--19.
\bibitem[Wolf(1992)]{Wolf1992} Wolf, S. (1992). Morality and partiality. \emph{Philosophical Perspectives}, 6, 243--259.
\end{thebibliography}

\end{document}
